\documentclass[11pt]{article}
\usepackage[margin=1in]{geometry}
\usepackage{amsthm,amsmath,amssymb,hyperref}
\newtheorem{theorem}{Theorem}
\title{Real zeros for a single lamination in a polygon}
\author{Wentao Li}
\date{November 20, 2026}
\begin{document}
\maketitle
\begin{abstract}
For any triangulation of an unpunctured marked disk, any arc and any single
lamination, the specialized cluster polynomial has only negative real zeros.
This proves the disk formulation of the log-concavity conjecture of Kang,
Lee and Lim. A single lamination divides the boundary vertices into two
intervals. Along the part of the triangulation crossed by both the arc and
the lamination, Ptolemy's relation reduces to two operations on a pair of
polynomials. Their ratio stays in an invariant complex sector.
\end{abstract}

\begin{theorem}
Let $T$ be a triangulation of a polygon, $L$ a single unweighted lamination
curve and $\gamma$ an arc. Specialize each initial mutable variable to one
and the frozen variable of $L$ to $q$. The resulting cluster variable
$c_\gamma(q)$ has positive constant term, nonnegative integer coefficients
and only negative real zeros. Its coefficient sequence is log-concave.
\end{theorem}

\begin{proof}
Draw the polygon convex. Up to isotopy $L$ is a chord between boundary-edge
interiors. Color vertices on one side $A$ and on the other side $B$, with the
starting vertex $p$ of $\gamma$ colored $A$. For any side or diagonal $ij$,
write $\delta_{ij}=1$ if its ends have different colors and $0$ otherwise.
This is precisely its minimal intersection number with $L$, also on the
boundary edges containing an end of $L$.

The tropical lambda length of $ij$ is $t_{ij}=q^{-\delta_{ij}/2}$.
By Definitions 14.6 and 15.3 and Theorem 15.6 of Fomin and Thurston
\cite{ft}, the specialized variable is the ordinary lambda length
$\ell_{ij}$ divided by $t_{ij}$. All edges of the initial triangulation,
including boundary edges, have $\ell_{ij}=t_{ij}$ after specialization.
Write $Y_v=\ell_{pv}/t_{pv}$. When $\gamma$ is in $T$ the answer is one.
Otherwise follow the triangles crossed by $\gamma$, starting with the triangle
at $p$. Its opposite edge is the first crossed diagonal, and both its endpoints
have $Y=1$.

At a crossed edge $ef$, let $w$ denote the third vertex of the next triangle.
The cyclic order of $p,e,w,f$ is as displayed, up to reversal. The edges
$ef,ew,fw$ all belong to the initial triangulation. Ordinary Ptolemy gives
\[
 Y_w = \frac{t_{pe}t_{fw}}{t_{pw}t_{ef}}Y_e+
        \frac{t_{pf}t_{ew}}{t_{pw}t_{ef}}Y_f.                       \tag{1}
\]
If $\gamma$ ends at $w$, return $Y_w$. Otherwise retain $w$ and whichever
of $e,f$ belongs to the next crossed edge.

The crossed edges with one end of each color occur consecutively. Indeed, the
dual graph of the triangulation is a tree. These edges are exactly the common
edges of the paths traversed by $\gamma$ and $L$, and two tree paths have
connected intersection. We can put the two curves in general position with
respect to the triangulation. When $e,f$ are both $A$, the two multipliers in
(1) equal one for either color of $w$. When both are $B$, the portion beyond
$ef$ is on the side of this chord away from $p$. Its vertices are all $B$:
the boundary arc between two $B$ vertices avoiding the interval of $A$ vertices
contains only $B$ vertices. Hence $w$ is $B$ and both multipliers again equal
one. The steps before the first mixed edge therefore give positive constants.
If there are no mixed edges, the answer is a positive constant.

At a mixed edge denote the values at its $A$ and $B$ ends by $U,V$,
respectively. Substitution in (1) gives
\[
       Y_w=U+V\quad(w\text{ of color }A),\qquad
       Y_w=qU+V\quad(w\text{ of color }B).                          \tag{2}
\]
While the next edge is still mixed, the update is either $(U+V,V)$ or
$(U,qU+V)$. All coefficients stay nonnegative and the constant terms stay
positive, including at the exit from the mixed portion.

Fix $z$ in the upper half-plane and let $\theta=\arg z\in(0,\pi)$.
At every mixed edge, $U(z)$ and $V(z)$ are nonzero and
$R=V(z)/U(z)$ lies in the closed sector $0\le\arg R\le\theta$.
This holds initially for positive constants. The second update sends $R$ to
$R+z$, which stays in the sector. The first sends it to $1/(1+1/R)$,
which stays in the sector because $1/R$ and $1+1/R$ have argument between
$-\theta$ and zero. All sums used here are nonzero since the sectors have
opening less than $\pi$.

If the calculation finishes at a mixed edge, both possibilities in (2) are
nonzero at $z$. If the next edge is monochromatic, its endpoint values are
either $U,U+V$ or $V,zU+V$. The ratios of new to retained value are $1+R$
and $1+z/R$, respectively, in the same sector. Subsequent steps use only
ordinary addition, so the final value is a nonnegative real linear combination
of this pair with at least one positive coefficient. It is nonzero. Thus
$c_\gamma$ has no zero in the upper half-plane. Real coefficients exclude the
lower half-plane and positivity excludes the nonnegative real axis. Newton's
inequalities give the log-concavity assertion.
\end{proof}

\section*{Two computations from the source}
Conjecture 1 of \cite{kll}, v4, is stated for an unpunctured disk. The broader
earlier surface formulation has annular obstructions \cite{tang,tilepe}.
The printed polynomial in Example 6 of v4 would violate log-concavity:
its three consecutive coefficients $10,6,4$ satisfy $6^2<10\cdot4$.
The triangulation drawn there gives a different polynomial. The source
of Figure 17 sets the number of vertices to twelve. With its labels,
start at vertex 6 and stop at vertex 11. The first crossed edge is $57$,
with vertex 7 of type $A$ and vertex 5 of type $B$. The subsequent third
vertices are $4,3,8,2,9,10,1,12,11$, of types $B,B,A,B,A,A,B,B,A$.
Formula (2) gives
\[
               c_{6,11}(q)=8q^3+26q^2+22q+5.
\]
In Example 4 the target crosses diagonals $1,2,3,4,5$ in that order.
Ordinary mutation in the displayed extended matrix gives the following
successive specialized variables:
\[
 1+q,\quad 2+q,\quad 3+2q,\quad 1+4q+2q^2,\quad 2+5q+2q^2.
\]
The last expression factors as $(2q+1)(q+2)$. An accompanying exact
seed-mutation script reproduces both examples from their triangulations.

\paragraph{Status and credit.} Research draft prepared with assistance from
OpenAI Codex. The conclusion covers the revised disk conjecture. Independent
proof and priority review is pending.
\begin{thebibliography}{9}
\bibitem{kll} W. Kang, K. Lee and E. Lim, \emph{Unimodality and Cluster
Algebras from Surfaces}, European Journal of Combinatorics 136 (2026),
104392. Disk formulation in the November 19, 2026 revision, Conjecture 1.
\url{https://arxiv.org/abs/2508.04396v4}
\bibitem{ft} S. Fomin and D. Thurston, \emph{Cluster algebras and
triangulated surfaces. Part II: Lambda lengths}, Memoirs of the AMS 255
(2018), no.~1223. \url{https://arxiv.org/abs/1210.5569}
\bibitem{tang} Q. Tang, \emph{Counterexamples to the Kang--Lee--Lim
log-concavity conjecture}. \url{https://aixiv.online/pdf/2611.00002}
\bibitem{tilepe} A. Til\`epe, \emph{Annular Obstructions and Skew-Shape
Block Identities}, 2026. \url{https://arxiv.org/abs/2611.14506v2}
\end{thebibliography}
\end{document}
